\section{Related Work}
\label{sec:related_work}

\subsection{Agent Trajectory Evaluation}
Unlike outcome benchmarks that inspect only terminal code patches or database records~\citep{jimenez2024swebench,qin2023toolllm,liu2023agentbench,zhou2024webarena}, trajectory evaluation examines intermediate planning decisions, tool invocations, and error recoveries generated during agent execution. Recent benchmarks expand assessment: BFCL~\citep{patil2025bfcl} evaluates multi-step tool calls via AST checks; TRAJECT-Bench~\citep{trajectbench2025} assesses tool selection and dependency order; AgentRewardBench~\citep{agentrewardbench2025} benchmarks trajectory judges; and DeepRed~\citep{deepred2026} and Claw-Eval~\citep{claweval2026} leverage environment snapshots to award partial credit along execution traces.

Despite providing intermediate visibility, existing evaluators face two key limitations. First, they predominantly rely on a single reference path~\citep{qian2024agentprocessbench}. In open environments, tasks frequently admit multiple valid solution routes; single-path matching inherently penalizes legitimate exploratory strategies. Second, they operate post-hoc across the full rollout and lack dynamic intervention to adapt evaluation strategies to an agent's specific weak links.

\subsection{Dynamic Evaluation for Agents}
Static benchmarks suffer from prompt overfitting, data contamination, and leaderboard score saturation. To address these vulnerabilities, researchers have begun exploring dynamic evaluation paradigms. For conversational agents, Agent-Testing Agent~\citep{agenttestingagent2026} adaptively generates subsequent test cases guided by prior judge feedback. In interactive tool environments, ToolSandbox~\citep{gu2024toolsandbox} introduces stateful conversational simulations with pre-defined milestones and safety minefields to evaluate arbitrary rollouts.

However, existing dynamic evaluation approaches primarily concentrate on dynamically selecting or perturbing benchmark test suites~\citep{gu2024toolsandbox,agenttestingagent2026}. During the actual execution of an individual task, trajectory evaluation remains largely passive: evaluators typically check simple milestone reachability at the end of rollouts, lacking a formal mechanism to dynamically steer evaluation focus, update dimension weights, or control execution progression in real time.

\subsection{Efficient Evaluation and Early Stopping Strategies}
As agent rollouts grow longer, evaluation costs have escalated dramatically, costing hundreds of dollars per benchmark pass~\citep{shi2025earlyeval}. Prior efficiency efforts centered on benchmark distillation, subsampling static suites into smaller representative subsets~\citep{aslam2006statistical,shi2025earlyeval}. To improve within-task efficiency, EarlyEval~\citep{shi2025earlyeval} introduces early outcome prediction, training dual LightGBM classifiers on prefix features to predict binary task success/failure and halt execution early.

Nevertheless, benchmark distillation reduces the task count while leaving long, wasteful rollouts untouched, whereas tabular early prediction models~\citep{shi2025earlyeval} act as black boxes. They predict binary outcomes without understanding domain semantics, cannot provide causal failure attribution, and lack awareness of tool dependency constraints or hazardous minefields.

In summary, existing trajectory evaluators either inspect isolated atomic steps that suffer from local noise, or evaluate full traces post-hoc against limited reference paths that penalize valid diversity.
Meanwhile, early stopping methods lack domain semantics and failure attribution.
To bridge this divide, we propose DynSTEER. DynSTEER performs stage-wise dynamic evaluation that adapts review capacity to weak or ambiguous stages. It compiles path-tolerant milestone graphs from available task inputs. Finally, it uses execution-time early stopping to eliminate wasteful computation.
